% !TITLE 诉衷情·当年万里觅封侯
% !AUTHOR 陆游
% !DYNASTY 两宋
% !GENRE 词
% !ORDER 16
% !EDITION 《全宋词》，唐圭璋编，中华书局
% !EDITION_STATUS pending

\begin{poem}
当年万里觅封侯\textsuperscript{1}，匹马戍梁州\textsuperscript{2}。
关河梦断何处\textsuperscript{3}？尘暗旧貂裘\textsuperscript{4}。

胡未灭\textsuperscript{5}，鬓先秋\textsuperscript{6}，泪空流。
此生谁料，心在天山\textsuperscript{7}，身老沧洲\textsuperscript{8}。
\end{poem}

\begin{annotations}
\notetext{1}{觅封侯：寻找建功立业、博取封侯的机会。用班超投笔从戎、立功西域被封为定远侯的典故。陆游四十八岁时曾到南郑（今陕西汉中）军中任职，是他一生中唯一一段真正的军旅生活。}
\notetext{2}{匹马戍梁州：单人独骑戍守梁州。梁州是古地名，这里指南郑一带，当时是南宋抗金的前线。}
\notetext{3}{关河梦断何处：关塞河防的军旅往事，如今只能在梦里追寻，醒来却不知身在何处。关河，指前线的关塞与河防。}
\notetext{4}{尘暗旧貂裘：从军时穿的貂皮袄积满灰尘，颜色都暗了。用《战国策》里苏秦的典故：苏秦游说秦王不成，黑貂之裘穿破，黄金用尽，只好落魄回家。这里说自己报国无门，白白老去。}
\notetext{5}{胡未灭：金人还没有消灭。胡，指北方的金国。}
\notetext{6}{鬓先秋：鬓发比心先老了。秋，秋天草木凋零，用来比喻人的衰老。}
\notetext{7}{天山：在今天的新疆，汉代是与匈奴作战的地方。这里借指南宋的抗金前线——他心心念念的仍然是战场。}
\notetext{8}{沧洲：水边的地方，指隐居之处。陆游晚年被罢官，闲居在山阴（今浙江绍兴）老家。“心在天山，身老沧洲”把心的去向和身的处境摆在一处，是这首词最沉痛的地方。}
\end{annotations}

\begin{appreciation}
《诉衷情》是陆游词里最沉痛的一首自述。他一生留下九千多首诗与一百多首词，篇幅大半都在反复申说同一件事：北伐。这首词只有四十四字，却把这件事从头说到了尾。

“当年万里觅封侯，匹马戍梁州”——当年为着建功立业奔走万里，单人独骑去戍守梁州。觅封侯用的是班超的典故，班超投笔从戎立功西域而受封定远侯。陆游四十八岁那年到南郑军中任职，前后不到一年，却是他一生中唯一一段真正的军旅生活。“万里”与“匹马”摆在一处，路越远，人越小，气势与孤独一同逼了出来。

“关河梦断何处？尘暗旧貂裘”——关塞河防的军旅往事如今只能在梦里追寻，一觉醒来却不知身在何处。貂裘用的是苏秦的典故，苏秦游说秦王不成，黑貂之裘穿破，黄金百斤用尽，只得落魄还家。陆游借这层意思说自己报国无门白白老去。一个“暗”字下得准，貂裘并没有破，只是积了灰，人被闲置得太久了。

“胡未灭，鬓先秋，泪空流”，三个三字句连成一串，一句逼一句，像三记闷拳，中间不给半点喘息的余地。金人还没有灭，鬓发倒先白了，眼泪白白地流。“先”字最狠，胡虏未灭而自己的头发倒先白了。这一串短句是整首词里最急促的一处。

“此生谁料，心在天山，身老沧洲”——谁料这一生，心还留在西北前线，人却要在江南水边老死。天山是汉代与匈奴作战的地方，这里借指南宋的抗金前线；沧洲是水边隐者所居，指他罢官后闲居的山阴老家。心上的地名与身上的地名遥遥对峙，一个在北，一个在南，中间隔着万里的山河与四十年的光阴。

南郑那不到一年的日子，陆游后来回忆了整整四十年，把那一年的大雪、射虎、铁衣与刁斗反反复复写进诗里。同一种错位辛弃疾也写过，“醉里挑灯看剑，梦回吹角连营”——醉里挑灯看的是剑，梦里回去的是营，睁眼却仍在江南。一个时代把两个最想打仗的人一并闲置在了水边。

“心在天山，身老沧洲。”词到这里便停住，天山在北，沧洲在南，八个字中间隔着一整个南宋的兴亡。
\end{appreciation}
